% !TeX root = ../sections/03_solutions.tex

\subsection{2_2}
\begin{solution}
	関数
	\[
	f(x)=|x|
	\]
	を
	\[
	T_n(x)
	=
	\frac{\alpha_0}{2}
	+
	\sum_{k=1}^{n}
	\left(
	\alpha_k\cos kx+\beta_k\sin kx
	\right)
	\]
	で近似する。
	
	平均 \(2\) 乗誤差を最小にする三角多項式は，
	\(f(x)\) のフーリエ部分和である。
	
	したがって，平均 \(2\) 乗誤差を最小にする係数は，
	\(f(x)\) のフーリエ係数である。
	
	つまり，
	\[
	\alpha_k=a_k,
	\qquad
	\beta_k=b_k
	\]
	である。
	
	この問題では \(\alpha_1\) と \(\beta_1\) を求めればよいので，
	\[
	\alpha_1=a_1,
	\qquad
	\beta_1=b_1
	\]
	を求めればよい。
	
	ここで，
	\[
	f(x)=|x|
	\]
	は偶関数である。
	したがって，正弦係数はすべて \(0\) である。
	
	よって，
	\[
	\beta_1=0
	\]
	である。
	
	次に，
	\[
	a_1
	=
	\frac{1}{\pi}
	\int_{-\pi}^{\pi}
	|x|\cos x\,dx
	\]
	を計算する。
	
	\(|x|\cos x\) は偶関数なので，
	\[
	a_1
	=
	\frac{2}{\pi}
	\int_0^\pi x\cos x\,dx
	\]
	である。
	
	部分積分を用いる。
	\[
	u=x,\qquad dv=\cos x\,dx
	\]
	とおくと，
	\[
	du=dx,\qquad v=\sin x
	\]
	である。
	
	したがって，
	\[
	\int_0^\pi x\cos x\,dx
	=
	[x\sin x]_0^\pi
	-
	\int_0^\pi \sin x\,dx.
	\]
	
	ここで，
	\[
	[x\sin x]_0^\pi=0
	\]
	であり，
	\[
	\int_0^\pi \sin x\,dx
	=
	[-\cos x]_0^\pi
	=
	-\cos\pi+\cos 0
	=
	1+1
	=
	2
	\]
	である。
	
	したがって，
	\[
	\int_0^\pi x\cos x\,dx=-2
	\]
	である。
	
	よって，
	\[
	a_1
	=
	\frac{2}{\pi}(-2)
	=
	-\frac{4}{\pi}.
	\]
	
	したがって，
	\[
	\boxed{
		\alpha_1=-\frac{4}{\pi},
		\qquad
		\beta_1=0
	}
	\]
	である。
	
	なお，問題文では \(n=30\) とされているが，
	\(\alpha_1,\beta_1\) の値は第 \(1\) フーリエ係数で決まるため，
	\(n=30\) であること自体はこの計算には影響しない。
\end{solution}